\subsection{维数与平坦性}

设 \(\rho: A \to B\) 是一个环同态。我们称下述条件为 (PM)：

(PM) 存在一个对 A 忠实平坦的 B-模  ，使得对于 B 的每个素理想  ，有 \(_{B}\)。

注。--- 1) 当存在一个有限生成的 B-模（它对 A 忠实平坦且支集等于 Spec(B)）时，条件 (PM) 得到满足。特别地，当 B 作为 A-模忠实平坦时，情况如此。

\begin{enumerate}
\def\labelenumi{\arabic{enumi})}
\setcounter{enumi}{1}
\item
  一个对 A 忠实平坦的 B-模  的存在性蕴含 \(\rho\) 的单射性（I, § 3, no. 5, 命题 8），以及映射 \({}^{a}\rho : \operatorname{Spec}(B) \to \operatorname{Spec}(A)\) 的满射性（II, § 2, no. 5, 命题 11 的推论 4）。
\item
假设 \(\rho: A \to B\) 是局部环之间的局部同态，并且存在一个 B-模 \(N\)，它在 A 上平坦，使得对 B 的每个素理想 \(A\)\(N \otimes_B \kappa(\mathfrak{q}) \neq 0\)\(B\)。那么 \(N\)\(A\) 在 A 上忠实平坦，因此 \(\rho\) 具有性质 (PM)：事实上，我们有  ，从而 \(N/\mathfrak{m}_{B}N = N \otimes_B \kappa(\mathfrak{m}_{B}) \neq 0\)，并且特别有 \(N \neq \mathfrak{m}_{B}N\)\(N \neq \mathfrak{m}_{A}N\)，结论由 I, § 3, no. 1 的命题 1 得出。
\end{enumerate}

命题 1. --- 设 \(\rho: A \to B\) 是满足条件 (PM) 的环同态。

\begin{enumerate}
\def\labelenumi{\alph{enumi})}
\item
  设 \(h: \mathbf{A} \to \mathbf{A}'\) 是一个环同态。这个同态 \(\rho': \mathbf{A}' \to \mathbf{A}' \otimes_{\mathbf{A}} \mathbf{B}\) 由 \(\rho\) 诱导，并满足条件 (PM)。
\item
  设  是 B 的一个素理想，\(^{-1}\)。典范同态 \(_{p}\)\(_{q}\) 满足条件 (PM)。
\item
设  是 B 的一个素理想，\(\mathfrak{p} = \rho^{-1}(\mathfrak{q})\)。对于包含于  中的 A 的每个素理想，存在 B 的一个素理想 \(\mathfrak{p}'\)，它位于 \(\mathfrak{q}'\) 之上且包含于 \(p'\)。
\end{enumerate}

设  是一个对 A 忠实平坦的 B-模，使得对于 B 的每个素理想 \(N \otimes_{B} \kappa(q) \neq 0\)。

证明 \(a)\)。A'-模 \(\mathbf{N}' = \mathbf{A}' \otimes_{\mathbf{A}} \mathbf{N}\) 是忠实平坦的（I, § 3, no. 3, 命题 5）；设  是 \(\mathbf{B}' = \mathbf{A}' \otimes_{\mathbf{A}} \mathbf{B}\) 的一个素理想， 是它在 B 中的逆像。我们有同构

\[
\mathrm{N} ^ {\prime} \otimes_ {\mathrm{B} ^ {\prime}} \kappa (\mathfrak {q} ^ {\prime}) \rightarrow \mathrm{N} \otimes_ {\mathrm{B}} \kappa (\mathfrak {q} ^ {\prime}) \rightarrow (\mathrm{N} \otimes_ {\mathrm{B}} \kappa (\mathfrak {q})) \otimes_ {\kappa (\mathfrak {q})} \kappa (\mathfrak {q} ^ {\prime});
\]

由于  ，我们还有 \(_{B}\)\(_{B'}\)。

证明  。根据 II, § 3, no. 4 的命题 13 和 14，\(_{p}\) 是平坦的。另一方面，设 \(_{q}\) 是 \(_{q}\) 的一个素理想；它具有形式  ，其中 \(_{q}\) 是包含于 \(\mathfrak{b}\) 的 B 的一个素理想（II, § 3, no. 1, 命题 3）。根据 N 的假设，我们有 \(\mathfrak{q}\)，又由于  同构于 \(_{q}\)，故有 \(_{q}\)\(_{q}\)。由注 3 可得结论。

证明 c)。由 \(\rho_{\mathfrak{q}}: \mathrm{A}_{\mathfrak{p}} \to \mathrm{B}_{\mathfrak{q}}\) 诱导的局部同态由 \(\rho\) 得到，并根据 b) 满足条件 (PM)。因此，映射 \(\operatorname{Spec}(\mathrm{B}_{\mathfrak{q}}) \to \operatorname{Spec}(\mathrm{A}_{\mathfrak{p}})\) 是满射（注记 2），这正是要证明的。

推论。--- 设 F 是 Spec(A) 的一个闭子集。如果 Y 是 F 在映射 \(^{a}\rho : \operatorname{Spec}(B) \to \operatorname{Spec}(A)\) 下的逆像的一个不可约分支，则 \(^{a}\rho(Y)\) 的闭包是 F 的一个不可约分支。

设  是 A 的一个理想，使得  ，并设  是 B 的素理想，使 \(F = V(a)\)。F 在  下的逆像是  的子集 \(Y = V(\mathfrak{q})\)，而  的闭包是不可约闭子集 \(V(\rho(a)B)\)，它包含于 \(V(\rho^{-1}(q))\)。

我们需证明：如果  是 B 的包含  的素理想中的极小元，则 \(\rho(a)\) 是 A 的包含 \(\rho^{-1}(q)\) 的素理想中的极小元。否则，A 中会存在一个素理想 \(p'\)，使得  且 \(a \subset p' \subset \rho^{-1}(q)\)；根据命题 1 c)，B 中会存在一个素理想 \(p' \neq \rho^{-1}(q)\)，使得 \(\mathfrak{q}'\) 且 ，从而 \(\mathfrak{q}' \subset \mathfrak{q}\) 且 \(\mathfrak{p}' = \rho^{-1}(\mathfrak{q}')\)，这与关于 \(\rho(a) \subset q' \subset q\)\(q' \neq q\) 的假设矛盾。

命题 2. --- 设  是具有性质 (PM) 的非零环同态。则有不等式

\[
\dim (\mathbf {B}) \geqslant \dim (\mathbf {A}) + \inf _ {\mathfrak {m} \in \mathbf {S}} \dim (\mathbf {B} / \mathfrak {m B})\tag{1}
\]

其中 S 是 A 的极大理想集合。

我们知道 \(\dim(A) = \sup_{\mathfrak m \in S} \dim(A_{\mathfrak m})\)\(\S\)\(n^o\)（§ 1, no. 3, 命题 8）。因此只需证明不等式

\[
\dim (\mathbf {B}) \geqslant \dim (\mathrm{A} _ {\mathfrak {m}}) + \dim (\mathrm{B} / \mathfrak {m B})\tag{2}
\]

对 A 的每个极大理想 m。换言之，需要证明不等式

\[
\dim (\mathbf {B}) \geqslant n + r\tag{3}
\]

如果 \(\mathfrak{p}_0\subset \ldots \subset \mathfrak{p}_n\) 是包含于  中的 A 的素理想链，且 \(\overline{\mathfrak{q}}_0\subset \ldots \subset \overline{\mathfrak{q}}_r\) 是  的素理想链，则对于 \(0\leqslant i\leqslant r\)，存在一个位于  之上的 B 的素理想，记为 \(\mathfrak{q}_{n+i}\)，并且 \(\overline{\mathfrak{q}}_i = \mathfrak{q}_{n + i} / \mathfrak{mB}\) 是 B 的素理想链。置 \(\mathfrak{q}_n\subset \ldots \subset \mathfrak{q}_{n + r}\)；令 \(\mathfrak{p}_i' = \mathfrak{p}_i\) 对 \(0\leqslant i\leqslant n-1\) 成立，从而 \(\mathfrak{p}_n' = \mathfrak m\)，且 \(\mathfrak{p}_0'\subset \ldots \subset \mathfrak{p}_n'\) 是 A 的素理想链，并且 \(\mathfrak{q}_n\) 位于 \(\mathfrak{p}_n'\) 之上。如果 \(\mathfrak{q}_i\) 是位于 \(\mathfrak{p}_i'\) 之上的 B 的素理想（\(1\leqslant i\leqslant n\)），则命题 1 c) 表明存在一个位于 \(\mathfrak{q}_{i-1}\) 和 \(\mathfrak{p}_{i-1}'\) 之上且包含于 \(\mathfrak{q}_i\) 的素理想。通过递降归纳可构造一条 \(\mathfrak{q}_0\subset \ldots \subset \mathfrak{q}_n\) 的 B 的素理想链，使得 \(\mathfrak{q}_i\) 位于 \(\mathfrak{p}_i'\) 之上，对 \(0\leqslant i\leqslant n\) 成立。由于 \(\mathfrak{q}_0\subset \ldots \subset \mathfrak{q}_{n+r}\) 是 B 的素理想链，我们证明了不等式 (3)。

注 4. --- 设 \(\rho: A \to B\) 是满足条件 (PM) 的诺特局部环之间的局部同态。我们将在稍后（§ 3, no. 4, 命题 7）看到，此时 (1) 中等式成立。在一般情形中不等式可能是严格的（p. 84, 练习 1）。

推论。--- 对 A 的每个理想  ，有
是 B 的一个包含  的素理想
，且 \(\rho(a)\mathbf{B}\)。根据命题 1，局部同态 \(p = \rho^{-1}(q)\)\(\rho_{\mathfrak{q}}: A_{\mathfrak{p}} \to B_{\mathfrak{q}}\) 由 \(\rho\) 诱导并满足 (PM)，因而由命题 2 有  。根据 § 1, no. 3 的命题 7，有  ，因为 \(\dim(A_{\mathfrak{p}}) \leqslant \dim(B_{\mathfrak{q}})\) 包含  ，从而对于包含 \(\leqslant \dim(A_p)\) 的 B 的每个素理想  ，有 \(\leqslant \dim(B_q)\)\(\rho(\mathfrak{a})\mathbf{B}\)。推论由同上得出。

引理 1. — 设 \(\rho: \mathbf{A} \to \mathbf{B}\) 是一个环同态，\(\mathfrak{p}\) 是 A 的素理想。与连续映射 \(^{a}h: \operatorname{Spec}(\mathbf{B} \otimes_{\mathbf{A}} \kappa(\mathfrak{p})) \to \operatorname{Spec}(\mathbf{B})\) 相联系的典范同态 \(h:\mathbf{B}\to\mathbf{B}\otimes_{\mathbf{A}}\kappa(\mathfrak{p})\) 把 \(\operatorname{Spec}(\mathbf{B}\otimes_{\mathbf{A}}\kappa(\mathfrak{p}))\) 同胚地映到由 \((^{a}\mathfrak{p})^{-1}(\mathfrak{p})\) 给出的 \(\operatorname{Spec}(\mathbf{B})\) 的子空间，该子空间由 \(\mathbf{B}\) 中位于 \(\mathfrak{p}\) 之上的素理想组成。

同态 \(h\) 由从 B 到 \(\mathbf{B} / \rho(\mathfrak{p})\) B 的商同态，以及从 \(\mathbf{B} / \rho(\mathfrak{p})\) B 到其分式环 \((\rho(\mathbf{A} - \mathfrak{p}))^{-1}(\mathbf{B} / \rho(\mathfrak{p}) \mathbf{B})\) 的典范同态复合而成。根据 II, § 4, no. 3 的命题 13 的注记和推论，\(^a h\) 因此把 \(\text{Spec}(\mathbf{B} \otimes_{\mathbf{A}} \kappa(\mathfrak{p}))\) 同胚地映到 \(\text{Spec}(\mathbf{B})\) 的子空间，该子空间由 B 的素理想 \(\mathfrak{q}\) 组成，这些素理想包含 \(\rho(\mathfrak{p})\) 且不交于 \(\rho(\mathbf{A} - \mathfrak{p})\)，即位于 \(\mathfrak{p}\) 之上。

注 5. --- 根据命题 2 和引理 1，在命题 2 的假设下，我们因此有不等式

\[
\dim (\operatorname{Spec} (B)) \geqslant \dim (\operatorname{Spec} (A)) + \inf _ {\mathfrak {p} \in \operatorname{Spec} (A)} \dim \left(({}^{a}\rho)^{-1} (\mathfrak {p})\right).\tag{4}
\]

